\section{哪些观测等价真的来自群作用}

\subsection{群的存在与有用群的存在不同}

所有满足 $F\circ T=F$ 的集合双射组成 $\operatorname{Aut}_{\rm set}(F)$。同一纤维内任意两点可以互换而固定其余点，所以在纯集合论上，任意等价关系都可以写成某个巨大群的轨道。\textbf{因此不能无条件声称一般纤维没有任何群作用。}

研究所需的对象应更严格：作用可计算、保持合法材料域、在声明的开放区域成立，且具有有限维、有限生成、光滑、代数或可证明的局部结构。对
\begin{equation}
\operatorname{Aut}_{\rm diff}(F)=\{T\in\operatorname{Diff}(\M):F\circ T=F\},
\end{equation}
封闭性来自 $F\circ T_1\circ T_2=F$，单位元显然存在，由 $F\circ T=F$ 右复合 $T^{-1}$ 得 $F\circ T^{-1}=F$，结合律来自函数复合。问题不是公理是否成立，而是能否找出非平凡元素、低成本执行，并证明轨道等于完整纤维。

任意优化正确器、最近点投影或饱和裁剪可能非单射，没有逆。将这些动作与几条弱方向拼起来，并不能得到一个群。它们可以构成有定义域的算法变换或对应关系，科学含义应按实际性质命名。

\begin{proposition}[精确光滑对称的必要条件]
若 $T$ 为局部微分同胚且 $F\circ T=F$，则
\begin{equation}
DF_{T(m)}DT_m=DF_m,
\end{equation}
从而
\begin{equation}
\rank DF_{T(m)}=\rank DF_m,\qquad
\im DF_{T(m)}=\im DF_m.
\end{equation}
同一光滑群轨道中的点必须有相同响应秩和响应切空间。
\end{proposition}
\begin{proof}
对观测保持恒等式应用链式法则；$DT_m$ 可逆，右乘不会改变线性映射的像与秩。
\end{proof}

这提供一个实际的排除测试：对已找到的精确材料逆解对，先比较材料有效 Jacobian 的像。若两个响应切空间明显不同，就不应再拟合一个在邻域成立的精确对称。像相同只是必要条件；高阶映射芽、合法域和全局拓扑仍可能不同。

\subsection{完整群轨道的充分条件与拓扑限制}

若存在材料区域的微分同胚 $\Phi:\M_0\to U\times G/K$，使 $F=h\circ\operatorname{pr}_U\circ\Phi$ 且 $h$ 单射，便可把 $G$ 在齐性空间 $G/K$ 上的作用拉回材料域，得到恰好等于全部纤维的轨道。这是明确的充分条件，难点在于实际构造 $\Phi$。更一般，一个已证明的主 $G$ 丛 $\M_0\to U$ 也有全局右作用和完整纤维轨道，不需要存在全局 section。

反过来，若一个 Lie 群完整作用且某纤维正好是一条轨道，该纤维应具有相应齐性空间 $G/G_m$ 的结构；连通群的轨道连通，紧群轨道紧，有限群轨道有限且大小由稳定子决定。这些性质都可以排除某些预设简单群。不同纤维可以有不同稳定子和轨道维数，故 rank 变化不排除一切群作用，但同一轨道内仍受前述必要条件约束。

常秩只给局部乘积。即便 proper submersion 给出局部平凡纤维丛，也没有自动给出一个指定有限维群的全局作用或廉价生成元；过渡结构群可能需要整个纤维的微分同胚群。正则局部 $\R^{d-r}$ 平移可在 chart 内使用，合法域截断后通常只形成局部伪群。

\subsection{解析正例与反例}

\begin{example}[$z^n$ 的完整循环轨道]
在 $\C$ 上取整数 $n\ge2$，$H(z)=z^n$，$T_kz=e^{2\pi ik/n}z$ 定义 $C_n$ 作用。非零纤维恰为 $n$ 个根，作用自由且传递；零点轨道只有一个点，稳定子为整个 $C_n$，实 Jacobian 在此退化。$\C\setminus\{0\}\to\C\setminus\{0\}$ 是 $n$ 层覆盖。

若限制实数材料，奇数次幂单射，偶数次幂只有 $\pm$ 两支；正材料域又只保留一支。复根的 $C_n$ 不能自动移植到被动实材料。非径向 prior 下，各根条件权重也不相等。
\end{example}

\begin{example}[圆与弯曲等值线都可以有连续群]
$H(a,b)=a^2+b^2$ 的正半径纤维是 $SO(2)$ 轨道，原点为固定点。另一映射 $H(a,b)=a+b^2$ 有全局 $\R$ 作用
\begin{equation}
T_t(a,b)=(a-2bt-t^2,b+t),\qquad T_tT_s=T_{t+s}.
\end{equation}
直接代入得 $H(T_t(a,b))=H(a,b)$，且任意同纤维两点可由某个 $t$ 连接。它在坐标 $(u,w)=(a+b^2,b)$ 中就是平移。因此纤维弯曲并不排除简单群；加入材料边界后，全局作用可能只剩局部或部分作用。
\end{example}

\begin{example}[三个根不意味着全局 $C_3$]
令 $f(x)=x^3-3x$。$f^{-1}(0)=\{-\sqrt3,0,\sqrt3\}$，但 $\R$ 上的有限阶保向 homeomorphism 必为恒等，反向 homeomorphism 若有限阶只能为二阶，故不存在有效全局 $C_3$。

事实上保持这个 $f$ 的全局 homeomorphism 只有恒等。$y>2$ 时逆像唯一，变换必须固定右侧整条半轴，故不能反向；保向变换在每个有限有序纤维上只能逐点固定。在临界值 $2$ 的纤维中，$-1$ 与 $2$ 的导数分别为 $0,9$，连光滑局部对称的必要秩条件都不满足。

限制到 $f^{-1}((-2,2))$ 的三个不相交区间后，可以用已知逆分支人为构造分支置换。这并没有免费找到逆分支，也不能跨临界值扩展为原全局对称。
\end{example}

\begin{example}[满秩的远处逆解也未必形成局部对称]
取
\begin{equation}
H(t)=(t^2-1,\ t(t^2-1)).
\end{equation}
$H(1)=H(-1)=(0,0)$，两点导数分别为 $(2,2)$ 与 $(-2,2)$，均满列秩，但响应切空间不同。因此不存在把两邻域互换且保持 $H$ 的局部微分同胚。这个映射处处秩为一；远处两个正则分支相交，并不自动构成稳定的多层覆盖。
\end{example}

有群也不代表轨道覆盖所有逆解。例如 $H(a,b)=a^2$，仅采用沿 $b$ 的平移群，会生成每条竖线，但当 $a\ne0$ 时同一观测纤维还包括另一条 $-a$ 竖线。需要再加入离散反射才能连接两支。即使能写成半直积，也须证明离散变换对连续变换的共轭作用、定义域和合法性；不能从两个模块的存在直接宣布半直积结构。

\subsection{指定群在 Maxwell 中的资格}

\begin{longtable}{p{.16\textwidth}p{.38\textwidth}p{.38\textwidth}}
\toprule 群或变换 & 合法出现的情形 & 固定采集下常见误判\\\midrule\endhead
$C_n,S_n$ & 真实几何自同构逐个固定源和接收，或显式可证的有限材料交换。 & 生成器中无序基元的标签交换往往只改变表示；同时重排观测标签通常只是等变。\\
$SO(2)$ & 连续旋转真正保持全部源场、接收泛函、边界、极化及材料域。 & 旋转环形阵列通常置换通道；一般有标签复数据并不不变。\\
$U(1)$ & 相位丢失观测或明确的未知相位 nuisance 模型可能含相关对称。 & 固定已知源与相干复数据下，统一转动电流相位一般改变观测，也未必满足构成方程。\\
$(\R,+)$ & 全局可平凡化纤维的平移，或精确垂直场的局部流。 & 接收零空间加法一般不保持共享材料二次约束。\\
$(\R_{>0},\times)$ & 无边界正参数的 $V=QD$ 是 decoder 尺度 gauge。 & 饱和、pure-cell 和属性边界破坏通用尺度商；几何与频率同时缩放改变了采集合同。\\
微分同胚群 & 同时允许变换介电率与磁导率张量的转换光学材料类。 & 固定磁导率和标量介电率的 P1 不包含这类轨道。\\\bottomrule
\end{longtable}

有限群和紧 Lie 群的不变量理论已有成熟成果。对有限维线性紧群作用，有限多项式不变量可以分离轨道，光滑不变量也可经这些不变量因子化；这里的前提是\emph{先给定真实作用}，不是通过数据拟合几个不变量就证明 Maxwell 有该群。[R3,R4] 一般非紧群可能有非闭轨道；连续不变量无法区分闭包相交的不同轨道，也不保证完整分类所有闭包，商可能不 Hausdorff。

\subsection{一个真正连续的 Maxwell 材料等价群}

允许三维 $\varepsilon,\mu$ 同时为各向异性张量。令 $\phi$ 是在源、接收及外部区域恒等的光滑保向微分同胚，$x'=\phi(x)$，$D=D\phi(x)$。变换
\begin{equation}
\varepsilon'(x')=\frac{D\varepsilon(x)D^\top}{\det D},\qquad
\mu'(x')=\frac{D\mu(x)D^\top}{\det D}
\end{equation}
配合 Maxwell 场的协变变换，保持外部带标签数据。这是转换光学和各向异性逆边界问题的已知非唯一性。光滑、非奇异变换可保持正定性与相应被动性；它不要求用奇异隐身斗篷。[R5]

这确实能改变固定物理坐标中的材料张量，故不同于电磁势 gauge。但原 P1 固定 $\mu=\mu_0I$。若变换后仍要求同一常数磁导率，必须有
\begin{equation}
DD^\top=(\det D)I.
\end{equation}
在三维取行列式且 $\det D>0$，得 $\det D=1$，进而 $DD^\top=I$。连通欧氏域中的光滑局部等距是刚性运动；外部恒等迫使整体恒等。因此这套已知连续群不能证明原标量、固定磁导率材料模型有对应自由度。

\section{群胚 分层结构与局部精确纤维}

\subsection{关系群胚的作用与限制}

对 $H:X\to Y$，$\mathcal R_H=X\times_YX$ 的箭头 $(x,x')$、单位 $(x,x)$、逆 $(x',x)$ 与首尾相接的复合确实定义等价关系群胚。该构造只记录谁和谁相等，没有告诉我们如何找到箭头。即使来自群作用，关系群胚也会忘记稳定子：作用群胚可有多个不同群元素连接同一对点，薄的关系群胚只有一条箭头。

若 $H:X^d\to U^r$ 是对共同观测流形的 submersion，则 $X\times_UX$ 是 $2d-r$ 维光滑流形，source 和 target 为 submersion，形成 Lie groupoid。\textbf{单独说 $H$ 全局常秩还不够}：上一节自交曲线的关系是对角线加两个孤立交叉箭头，这些箭头的 source 不能成为到一维基的 submersion。

若 $F$ 为有限维实半代数或解析映射，可在相应条件下使用分层。材料合法边界、field-node、判别式和 rank 变化分别划分 strata。一般光滑 learned decoder 不自动给有限半代数或 Whitney 分层；需要另加正则性条件或只声明实际构建的局部 chart。奇异叶状结构的 holonomy groupoid 也需要指定局部有限生成、对 Lie bracket 封闭的向量场模，不由一句 ``取弱子空间'' 得到。[R6]

\subsection{精确 kernel 的可积性}

\begin{theorem}[正则局部纤维]
设 $H:U\subset\R^d\to\R^p$ 为 $C^k$，$k\ge1$，且在 $x_0$ 邻域常秩 $r$。则经过 $x_0$ 的局部等观测集为 $d-r$ 维嵌入流形，其切空间为 $\ker DH$。若 $r=d$，$H$ 在 $x_0$ 的某邻域单射。
\end{theorem}
\begin{proof}
常秩定理将映射写为 $(u,v)\mapsto(u,0)$，固定 $u$ 后 $v$ 可自由变化。满列秩时选取一个可逆的 $d\times d$ 输出 minor，对相应输出坐标应用逆函数定理。[R1]
\end{proof}

若光滑向量场 $V,W$ 始终满足 $DH\,V=DH\,W=0$，对每个分量 $H_j$，
\begin{equation}
[V,W]H_j=V(WH_j)-W(VH_j)=0.
\end{equation}
因此 exact vertical vector fields 的 bracket 仍精确垂直。常秩内其积分叶是局部纤维的连通片；全局一条纤维可含多个连通分量，不能由一条局部叶代表。

\subsection{弱方向的非封闭性}

\begin{example}[小泄漏不能控制 bracket 泄漏]
令 $H(x,y,z)=z$，
\begin{equation}
V=\partial_x,\qquad W=\partial_y+\epsilon\sin(x/\epsilon)\partial_z.
\end{equation}
$|VH|=0$、$|WH|\le\epsilon$，但
\begin{equation}
[V,W]=\cos(x/\epsilon)\partial_z,
\end{equation}
在 $x=0$ 处的观测泄漏为一。小的有限噪声响应并不控制方向场的空间变化率。
\end{example}

真正的截断 SVD 弱分布也可非可积。例如 $H_\epsilon(x,y,z)=(\epsilon x,\epsilon y,z+xy)$，$0<\epsilon<1$。令 $\tau=\epsilon^2$、$s=x^2+y^2$、
\begin{equation}
\lambda=\frac{\tau+s+1+\sqrt{(\tau+s+1)^2-4\tau}}2,
\qquad \kappa=\frac{\lambda}{\lambda-\tau}.
\end{equation}
$DH^\top DH$ 的最强特征向量为 $(\kappa y,\kappa x,1)$，另两特征值为 $\tau,\tau/\lambda\le\tau$。弱平面由 $\alpha=dz+\kappa(s)(y\,dx+x\,dy)$ 消去，而
\begin{equation}
\alpha\wedge d\alpha=2\kappa'(s)(x^2-y^2)\,dx\wedge dy\wedge dz
\end{equation}
一般非零，故不满足 Frobenius 条件。实验中的 projector bracket 只能说明选定谱截断分布的性质；它不是 exact fiber 的拓扑证据。

R7 八个中心的有效响应数值秩 66，最小材料响应奇异值约 $0.0308$--$0.0956$，所选单位材料速度的一阶白化响应约 $8.33$--$16.53$。这些是满秩但有限噪声较弱的材料方向，不是 $DHv=0$。若数学上确为满列秩，光滑观测保持 Lie 群的恒等连通分支在该局部区域只能平凡作用；远处离散作用仍可能存在。[I2]

\subsection{高阶补偿何时具有组合结构}

若 $H$ 常秩且允许所有合法切向补偿，每个 exact-kernel 速度都能提升为精确等观测曲线；有限阶障碍因此可消去。若只有某条固定直线的二阶、三阶项消失，则远不足以证明曲线存在。例如 $H(t)=t^{k+1}$ 的前 $k$ 阶在零点消失，而零观测的逆像仍孤立。

局部微分同胚的有限阶 jets 本身可以按截断复合组成 jet 群；不同基点对应 jet 群胚。保持给定 $H$ 的有限阶 jet 的集合也可以形成相应稳定结构，但必须限定相同基点、相容阶数和可逆的一阶项。它保证的是模掉高阶余项后的代数，不是有限步长上的物理观测相等。普通 GN corrector 经常是非可逆的，并不自动属于这类群。

若已经证明 $T_g$ 为一个作用，但只有近似观测保持，记
\begin{equation}
\delta_g(m)=W_a\{F(T_gm)-F(m)\}.
\end{equation}
则
\begin{equation}
\delta_{gh}(m)=\delta_g(T_hm)+\delta_h(m).
\end{equation}
这是误差复合的准确记账。每步小误差可以累积，不能把它改名为精确 cocycle 不变量或物理 holonomy。对 finite-noise 路径应始终回到同一观测锚点验收。

\subsection{扩展材料类后的严格正面维数结果}

旧 Research Z 已提出全尺度小波材料 prior，在 $24\times24$、固定复相位下保留 576 个实系数。[I3] 若完整基变换可逆、全部 576 个系数方差为正，且整个材料 decoder 的 Jacobian 在内点满列秩 576（例如逐 cell link 的每个导数均非零），合法材料像局部为 576 维。原采集的实数据维数为 $2\times8\times32=512$，故任何常秩区域满足
\begin{equation}
\dim\{m:F(m)=F(m_0)\}_{\rm local}=576-r\ge64.
\end{equation}
若进一步证实 $r=512$，满行秩在邻域保持，隐函数定理直接给出 64 维精确局部纤维。

这是新 prior $P_Z$ 下的条件定理。单个奇异点不能仅凭维数计数断言纤维维数，例如 $F(a,b)=a^2+b^2$ 在零观测仍只有原点。边界也可能截断纤维。即使正维纤维存在，也未证明可达幅度大、包含新拓扑或有显著 posterior mass；这些是后续任务与概率问题。

\section{全局分支 覆盖与 Monodromy}

\subsection{局部流能够覆盖什么}

\begin{proposition}[生成元的覆盖资格]
有限个固定变换及其逆的有限复合只能生成可数个点，不能精确列举正维连续纤维。若改用有限个连续参数光滑流，限定在同一合法、无边界的连通正则纤维 stratum 内，允许正负时间且所有流段不离开该域，并且这些精确垂直向量场处处张成全部切空间，则有限段流的可达轨道覆盖该连通纤维。若以 Lie closure 张成替代直接张成，应满足相应轨道定理的条件。
\end{proposition}
\begin{proof}
有限字母表的有限字词可数。连续流情形，各点局部流复合给出纤维中的开邻域；可达轨道将纤维分成互不交的开集，连通性迫使只有一个。一般括号生成情形使用 Sussmann 轨道定理。[R2]
\end{proof}

完整性、材料边界、节点和 blow-up 决定流可否继续。该结论不覆盖其他连通分量，也不允许精确可逆流把不同 rank 的 strata 随意连接。稠密子群可以逼近轨道，却不等于精确覆盖，更不提供原先验均匀性。

\subsection{有限覆盖的条件}

对参数化系统 $P(m;u)=0$，若在去除判别式和不适定点的基空间 $U^\circ$ 上，投影是有限无分歧覆盖，沿闭合 $u(t)$ 可唯一提升各逆解。回到同一 $u_0$ 后得到有限纤维的置换。这是 monodromy。若总覆盖空间连通，monodromy 作用传递；不同不可约或不连通分量可能需要不同种子。[R7]

须分别区别以下三个对象：固定 $u$ 上的分支集合；整个覆盖的 deck transformation；基空间闭路引起的 monodromy permutation。后者可能比 deck group 大得多，甚至在 deck group 平凡时仍能遍历全部分支。

\begin{example}[复三次映射的 monodromy]
$f(z)=z^3-3z$ 的复临界值为 $\pm2$。绕这两个点的闭路分别交换不同的两个根，生成 $S_3$；这个连通三层覆盖的 deck group 却平凡。本轮以 $y=0$ 的三根为起点，对两个闭路各作 800 步延拓，得到置换 $[1,0,2]$ 与 $[0,2,1]$，最大多项式残差约 $9.69\times10^{-14}$，终点匹配误差约 $1.63\times10^{-16}$。

这是复 toy 的实际数值核验。实基区间 $(-2,2)$ 无这种非平凡闭路；真实有序根不能不经碰撞而交换顺序。它不能被描述为被动物理材料路径。
\end{example}

对 Maxwell，测量空间往往比材料维数高。任意在 ambient $\Y_a$ 画一个圆通常离开物理观测像，逆问题根本无解。必须在实际观测像、可实现控制族或明确的复辅助多项式族内延拓。

复路径仍可以作为\emph{计算工具}发现最终为实、被动且属于 prior 支持的材料端点。路径不物理并不使端点无效，但它证明的是代数搜索到达，不是原实材料域中存在物理等响应路径。

\subsection{一个必须排除的伪循环：沿线段原路返回}

\begin{proposition}[往返路径不产生正则 monodromy]
若参数路径是某条路径 $\alpha$ 紧接其逆路径 $\alpha^{-1}$，且每段提升存在并满足唯一连续提升，则其作用于起始纤维的置换是恒等。特别是辅助系统
\begin{equation}
P(m)+\gamma\sin(2\pi t)=0
\end{equation}
的参数只沿同一线段来回运动。若全程保持正则、有限且可唯一跟踪，则 $t=1/2$ 和 $t=1$ 都返回原解，而不会自动产生另一支。
\end{proposition}
\begin{proof}
第一段提升反向后就是第二段的一个提升，唯一性迫使第二段沿原路返回。正弦路径前半周期先从零到 $-\gamma$ 再回零，后半周期对 $+\gamma$ 重复同样论证。
\end{proof}

若这样的数值过程返回不同根，应查找沿途 Jacobian 退化、无穷远、提升不唯一或 Newton 跳支。新根仍可能是合法候选，但不能据此宣称已经证明一个覆盖空间的非平凡 monodromy 或循环群。真正的覆盖 monodromy 必须使用在去判别式基空间中不能按这种方式消去的闭路，例如本轮三次多项式的复绕圈。

这一限制并不否定经过奇点的合法换根：对 $P(x)=x^2-1$、$\gamma=1$，$x(t)=\cos(\pi t)-\sin(\pi t)$ 满足上述辅助系统并从 $1$ 到 $-1$；但在 $t=1/4$ 穿过 $x=0$、$DP=0$。成功端点、奇异分支选择和正则覆盖的 monodromy 应分别记录。

\subsection{分支覆盖的五级记录}

\begin{enumerate}
\item \textbf{端点资格}：若干候选通过完整物理检查，根数只是下界。
\item \textbf{轨道资格}：已证明某个指定作用的全部不同轨道元素被表示。
\item \textbf{连通纤维资格}：条件定理或可认证 chart 覆盖给出某一连通正则分量。
\item \textbf{限定域全分支资格}：在预声明有界材料域内，全部可行根或分量都有覆盖或排除证书。
\item \textbf{全 prior 资格}：另有域外概率界、全部 strata 和原子处理，足以支撑全目标 posterior 的声明。
\end{enumerate}

多次 multi-start 或 monodromy 未发现新根，不会自动升级完整性。数值代数几何中的 witness set 描述复代数组件；一个复不可约组件的合法实部分仍可能有多个连通分量。单根的 $\alpha$-理论认证保证某个近似点收敛到唯一根，不证明其他根不存在。HomotopyContinuation 的官方文档也明确将数值 completeness verification 与严格认证区分。[R8,R9]
